% =====================================================================
%  setup.tex — الحزم والخطوط وتخطيط الصفحة
%  يُترجَم بـ XeLaTeX فقط (في Overleaf: Menu > Compiler > XeLaTeX)
% =====================================================================

% ---------- تخطيط الصفحة ----------
\usepackage[a4paper,top=2.4cm,bottom=2.4cm,inner=2.2cm,outer=2.2cm,headheight=16pt]{geometry}

% ---------- حزم عامة (تُحمَّل قبل polyglossia/bidi) ----------
\usepackage{fontspec}
\usepackage{xcolor}
\usepackage{graphicx}
\usepackage{array}
\usepackage{tabularx}
\usepackage{booktabs}
\usepackage{multirow}
\usepackage{longtable}
\usepackage{enumitem}
\usepackage{fancyhdr}
\usepackage{titlesec}
\usepackage{titletoc}
\usepackage[most]{tcolorbox}
\usepackage{etoolbox}
\usepackage{needspace}
\usepackage[hidelinks,bookmarksnumbered,unicode]{hyperref}

% ---------- اللغة: العربية لغة رئيسية ----------
\usepackage{polyglossia}
\setmainlanguage[numerals=mashriq]{arabic}
\setotherlanguage{english}

% ---------- الخطوط (مضمّنة في مجلد fonts/ لتعمل على أي جهاز وعلى Overleaf) ----------
\defaultfontfeatures{Ligatures=TeX}
\setmainfont{Amiri}[
  Path=fonts/,
  Extension=.ttf,
  UprightFont=*-Regular,
  BoldFont=*-Bold,
  ItalicFont=*-Italic,
  BoldItalicFont=*-BoldItalic,
  Script=Arabic,
  Scale=1.0]
\newfontfamily\arabicfont{Amiri}[
  Path=fonts/,
  Extension=.ttf,
  UprightFont=*-Regular,
  BoldFont=*-Bold,
  ItalicFont=*-Italic,
  BoldItalicFont=*-BoldItalic,
  Script=Arabic]
\newfontfamily\quranfont{AmiriQuran}[Path=fonts/,Extension=.ttf,Script=Arabic,Scale=1.05]
\newfontfamily\englishfont{Amiri}[Path=fonts/,Extension=.ttf,UprightFont=*-Regular,BoldFont=*-Bold,ItalicFont=*-Italic]

% ---------- المسافات ----------
\linespread{1.25}
\setlength{\parindent}{0pt}
\setlength{\parskip}{0.45em}
\setlist{itemsep=2pt,topsep=3pt,leftmargin=*}

% ---------- رأس الصفحة وتذييلها ----------
\pagestyle{fancy}
\fancyhf{}
\fancyhead[C]{\small\color{headcolor}\leftmark}
\fancyfoot[C]{\color{headcolor}\thepage}
\renewcommand{\headrulewidth}{0.4pt}
\renewcommand{\headrule}{\color{gold}\hrule height 0.6pt}
\renewcommand{\chaptermark}[1]{\markboth{#1}{}}
\fancypagestyle{plain}{\fancyhf{}\fancyfoot[C]{\color{headcolor}\thepage}\renewcommand{\headrulewidth}{0pt}}

% ---------- العناوين ----------
% الباب (chapter) بلا رقم ظاهر، بإطار مزخرف
\titleformat{\chapter}[block]
  {\centering\normalfont}
  {}{0pt}
  {\begin{tcolorbox}[enhanced,colback=parchment,colframe=gold,boxrule=1.2pt,
     arc=3mm,width=0.8\linewidth,center,halign=center,
     borderline={0.4pt}{3pt}{gold}]
     \Huge\bfseries\color{maroon}}
  [\end{tcolorbox}]
\titlespacing*{\chapter}{0pt}{-10pt}{18pt}

\titleformat{\section}[block]
  {\needspace{8\baselineskip}\normalfont\Large\bfseries\color{maroon}}
  {}{0pt}
  {}
  [{\color{gold}\titlerule[0.8pt]}]
\titlespacing*{\section}{0pt}{14pt}{8pt}

\titleformat{\subsection}[block]
  {\normalfont\large\bfseries\color{headcolor}}
  {}{0pt}{}
\titlespacing*{\subsection}{0pt}{10pt}{4pt}

\setcounter{secnumdepth}{-1}   % بلا ترقيم للعناوين
\setcounter{tocdepth}{1}       % الفهرس: الأبواب والأقسام

% تنسيق الفهرس
\titlecontents{chapter}[0pt]{\addvspace{6pt}\large\bfseries\color{maroon}}
  {}{}{\titlerule*[0.6pc]{.}\contentspage}
\titlecontents{section}[1.5em]{\normalsize}
  {}{}{\titlerule*[0.6pc]{.}\contentspage}
